% Appendix body only — the builder emits \section{...}\label{sec:app:gray}.
% Primary source: docs/derivations/G2c_gray_a9_a10_mapping.md (commission task G2c)
This appendix records the term-by-term audit of the likelihood
assembled in Section~\ref{sec:framework} against the published
equations of \citet{Gray:2019ksv}, whose appendix equations we cite as
A9, A10, etc.  The comparison was carried out against the equations as
published, not against secondary descriptions of them, and its outcome
is summarized in Table~\ref{tab:gray:mapping}: the assembled per-event
likelihood is an exact algebraic rearrangement of the published
hypothesis mixture, the audit found no algebraic discrepancy, and
every remaining difference is a deliberate, documented deviation ---
the most consequential being the localization-ball selection
denominator of \texttt{local\_ratio}.
% src: docs/derivations/G2c_gray_a9_a10_mapping.md (VERDICT)

\subsection{Exact correspondences}
\label{sec:app:gray:exact}

The hypothesis mixture of equation~\eqref{eq:mixture} is Gray
et~al.'s eq.~9.  Their eq.~A14 defines the catalogue-membership weight
$p(G \mid D_i, h)$ as a ratio of selection integrals; pre-integrating
their hard apparent-magnitude threshold (Schechter-function magnitude
and mass integrals) turns it into the smooth incompleteness weight
$1 - f(z,\Omega)$, applied per HEALPix pixel
\citep{MISSING:Gray2022-pixelated}.
% MISSING CITATION: R. Gray, C. Messenger, J. Veitch, "A pixelated approach to galaxy catalogue incompleteness: improving the dark siren measurement of the Hubble constant", MNRAS 512 (2022) 1127, arXiv:2111.04629
Under that rewriting, the numerator of eq.~A14 is
$\beta_G(h)$ [equation~\eqref{eq:betag}] and its denominator is
$D(h)$ [equation~\eqref{eq:Dh}]; eq.~A15 is the complement
$1 - w_G$; and the out-of-catalogue numerator and selection integrals,
eqs~A17 and A18, are $B_\mathrm{num}(h)$
[equation~\eqref{eq:bnum}] and $\beta_{\bar G}(h)$
[equation~\eqref{eq:betabar}], whose ratio is the completion-term
likelihood of eq.~A19.  Substituting all five identifications into
the mixture cancels the factors of $\beta_{\bar G}$ and yields the
single ratio of equation~\eqref{eq:assembled} --- an identical
rearrangement, not an approximation, and one that remains finite in
the complete-catalogue limit $f \to 1$, where eq.~A19 becomes $0/0$.
% src: docs/derivations/G2c_gray_a9_a10_mapping.md §1
Throughout, the selection rule of \citet{Mandel:2018mve} is preserved:
$p_\mathrm{det}$ enters only the denominators $D_g$, $D$,
$\beta_{\bar G}$ (and the global catalogue sum discussed below), never
a measurement numerator.

\begin{table}
  \caption{Term-by-term mapping of the implemented estimator onto
  \citet{Gray:2019ksv}.  ``Exact'' means an algebraic identity up to
  the smooth-completeness rewriting of the magnitude threshold;
  ``---'' means no published counterpart exists.}
  \label{tab:gray:mapping}
  \begin{tabular}{lll}
    \hline
    This work & \citet{Gray:2019ksv} & Status \\
    \hline
    $p_i(h)$, eq.~\eqref{eq:assembled} & eq.~9 & exact rearrangement \\
    $w_G = \beta_G/D$, eq.~\eqref{eq:betag} & eqs~A14--A15 & exact; smooth $f$ \\
    $D(h)$, eq.~\eqref{eq:Dh} & denom.\ of eq.~A14 & exact \\
    $\beta_{\bar G}(h)$, eq.~\eqref{eq:betabar} & eq.~A18 & threshold $\to$ smooth $1-f$ \\
    $B_\mathrm{num}(h)$, eq.~\eqref{eq:bnum} & eq.~A17 & exact sky marginal \\
    $B_\mathrm{num}/\beta_{\bar G}$ & eq.~A19 & implied; never divided \\
    $L_\mathrm{cat}$, eq.~\eqref{eq:est:lcat} & eqs~A9--A10 & denom.\ set deviates (cf.~A20) \\
    $N_g$, $D_g$, eqs~\eqref{eq:est:Ng}--\eqref{eq:est:Dg} & integrands of eq.~A10 & exact \\
    bare $p_g(z)$ [$\Pi \equiv 1$, eq.~\eqref{eq:kernel}] & $p(z_i)$ in eq.~A10 & literal \\
    deconvolved $p_g(z)$, eq.~\eqref{eq:est:pgz} & --- & Appendix~\ref{sec:app:voldeconv} \\
    $w_g = R_\mathrm{eff}(M_g)/(1+z_g)$ & $p(s\mid z)\,p(s\mid M)$; eqs~A3--A4 & instantiated \\
    MBH-mass channel & --- & Appendix~\ref{sec:app:eddm} \\
    \hline
  \end{tabular}
\end{table}

\subsection{The local selection denominator: nearest analogue eq.~A20}
\label{sec:app:gray:local}

In the notation of Section~\ref{sec:framework}, the in-catalogue
likelihood of Gray et~al.\ with per-galaxy redshift uncertainty, their
eq.~A10, reads
\begin{equation}
  p(x_i \mid G, D_i, h)
  \;=\;
  \frac{\displaystyle\sum_{g=1}^{N_\mathrm{gal}} \int \mathrm{d}z\;
        p\big(x_i \mid z, \Omega_g, h\big)\, p(s \mid z)\,
        p(s \mid M_g)\; p_g(z)}
       {\displaystyle\sum_{g=1}^{N_\mathrm{gal}} \int \mathrm{d}z\;
        p\big(D_i \mid z, \Omega_g, h\big)\, p(s \mid z)\,
        p(s \mid M_g)\; p_g(z)},
  \label{eq:gray:a10}
\end{equation}
% src: docs/derivations/G2c_gray_a9_a10_mapping.md §0 (A10 transcription)
where $p(s\mid z)$ and $p(s\mid M)$ are their generic host-probability
weights (their masses derive from apparent magnitudes,
$M(z_i, m_i, H_0)$; here $M_g$ is the catalogued MBH mass) and
\emph{both} sums run over the entire catalogue of $N_\mathrm{gal}$
galaxies --- no restriction to a localization region appears in the
published equation.  Restricting the numerator sum to the ball
$\mathcal{B}$ is numerically inert, since the measurement Gaussian
suppresses galaxies outside it exponentially; the faithful discrete
transcription of eq.~\eqref{eq:gray:a10} is therefore the
\emph{global} variant of Section~\ref{sec:pitfall}, which keeps the
denominator summed over the full catalogue.  Multiplying that discrete
sum by the continuous $\beta_G(h)$ in equation~\eqref{eq:assembled}
implicitly requires the two to agree in $h$-dependence up to a
constant; on GLADE+ this fails at the $-17$ per cent level end-to-end
across the posterior grid (Appendix~\ref{sec:app:betag}),
% src: .planning/gate/G1_beta_g_check.json (end_to_end_tilt_h3_corrected = -0.172)
and the discrete denominator additionally evaluates $p_\mathrm{det}$
at the catalogued $z_g$ while the numerator integrates over the
photometric kernel --- an asymmetric treatment of the redshift
smearing that is precisely the prior inconsistency of
Section~\ref{sec:pitfall}.
% src: docs/derivations/G2c_gray_a9_a10_mapping.md §6 (C6) and §4.1

The \texttt{local\_ratio} estimator [equation~\eqref{eq:est:lcat}]
instead evaluates both sums over the same ball with the same weights
and the same kernel.  This is \emph{not} eq.~A10; its nearest
published analogue is Gray et~al.'s catalogue-patch case, their
eq.~A20, in which the in-catalogue term is the regular ratio of sums
with the sky integrals limited to a patch.  The deviation identifies
the patch event-by-event with the three-dimensional localization ball
(sky \emph{and} distance) rather than a fixed RA/Dec.\ patch, and it
omits eq.~A20's explicit term for the sky remainder
$\Omega_\mathrm{rest}$, whose role --- hosts outside the patch --- is
played approximately by the completion channel, since $D(h)$ and
$\beta_{\bar G}(h)$ are built over the full sky and volume.
% src: docs/derivations/G2c_gray_a9_a10_mapping.md §4.2, §5 (D7)
The deviation is principled on three grounds: within the ball every
per-galaxy constant (number-density calibration, rate normalization)
cancels row by row, so no bridge between the discrete catalogue and a
continuous density model is invoked anywhere; numerator and
denominator treat the photometric kernel identically, restoring the
symmetry that the global sum breaks; and the resulting configuration
passes the coverage test of Section~\ref{sec:coverage}, where the
literal global transcription yields approximately zero empirical
coverage on photometric-redshift catalogues.
% src: docs/derivations/G2c_gray_a9_a10_mapping.md §4.1 (~0% P–P coverage of "global")

\subsection{Constructs beyond the published equations}
\label{sec:app:gray:beyond}

Four ingredients map to no equation in \citet{Gray:2019ksv} and are
cited here as project constructions.  (i)~The volume-deconvolved
kernel of equation~\eqref{eq:est:pgz} with its per-galaxy
normalization $Z_g$: eq.~A10 treats $p(z_i)$ as externally given, and
the mock catalogues of \citet{Gray:2019ksv} had zero redshift error;
the likelihood-times-prior construction parallels the line-of-sight
redshift priors of later \textsc{gwcosmo} releases
\citep{MISSING:Gray2023-los-zprior}.
% MISSING CITATION: Gray et al. 2023, "Joint cosmological and gravitational-wave population inference using dark sirens and galaxy catalogues", arXiv:2308.02281
(ii)~The local selection denominator, as above.  (iii)~The MBH-mass
channel in its entirety (Appendix~\ref{sec:app:eddm}): Gray et~al.'s
masses enter only through population priors in the selection
normalization, never as a per-galaxy mass-likelihood channel.
(iv)~The factorized product
$[1 - f_{k(\Omega_i)}(z)] \times (\sin\hat\theta_i/4\pi)\,
\mathcal{N}(u; 1, \Sigma_{i,uu})$ inside equation~\eqref{eq:bnum}:
each factor is published [eq.~A17 and the pixelated completeness of
\citealt{MISSING:Gray2022-pixelated}], but eq.~A17
strictly requires a single sky integral weighting both factors; the
factorization is exact only where $f$ is constant across the GW sky
support, an excellent approximation at the median localization of
$0.2\,\mathrm{deg}^2$, far below the completeness pixel scale.
% src: docs/derivations/G2a_completion_sky_marginal_4pi.md §4 (median 0.2 deg²); docs/derivations/G2c_gray_a9_a10_mapping.md §5 (D5)

One approximation deserves explicit statement: the rate weight $w_g$
is evaluated at the catalogued $(M_g, z_g)$ and held outside the
redshift integrals of equation~\eqref{eq:gray:a10}, whereas eq.~A10
carries $p(s \mid z_i)$ inside $\int \mathrm{d}z_i$.  For narrow
kernels the difference is second order; at GLADE+ photometric widths
it is not guaranteed to be, and we state it here as a known
approximation of the implementation rather than silently absorb it.
% src: docs/derivations/G2c_gray_a9_a10_mapping.md §5 (D1 sub-deviation)
